\chapter{Descripción del modelo experimental y simulado}
\label{cap:modelo}

Este capítulo describe con qué se obtuvieron los resultados del
capítulo~\ref{cap:resultados}: el modelo simulado, el montaje experimental y
los métodos de medición. Todas las herramientas son parte del repositorio y
pueden volver a ejecutarse.

\section{Modelo simulado}
\label{sec:modelo-simulado}

\paragraph{Simulación del núcleo completo}
El banco de prueba de Icarus Verilog, heredado del proyecto original y
adaptado al port, instancia el núcleo con un modelo de comportamiento de la
memoria externa, le entrega un flujo elemental MPEG-2 desde un archivo y
vuelca el contenido de la memoria de cuadros. El volcado se hizo idéntico,
palabra por palabra, al que se obtiene leyendo la DDR en la placa, de modo
que simulación y hardware se analizan con exactamente las mismas
herramientas.

\paragraph{Bancos de bloque}
Cada bloque nuevo tiene su banco propio (sección~\ref{sec:verificacion-rtl}):
puente APB, puente de memoria contra un esclavo AXI con latencias y
reordenamiento variables, DMA del flujo y ruta de respuestas de memoria con
el COREFIFO de Microchip.

\paragraph{Banco IEEE 1180}
Un banco dedicado alimenta la IDCT del núcleo, el mismo archivo Verilog que
se sintetiza, con los bloques pseudoaleatorios que define el estándar
IEEE~1180-1990 para sus seis condiciones de prueba (tres rangos de valores,
con los dos signos, 10\,000 bloques cada una) y compara la salida contra una
IDCT de referencia en doble precisión.

\section{Modelo experimental}
\label{sec:modelo-experimental}

\paragraph{Montaje}
La plataforma es una PolarFire SoC Discovery Kit (dispositivo MPFS095T,
encapsulado FCSG325), alojada en una caja estanca para protegerla durante el
desarrollo. Se conecta a una computadora de escritorio con Linux por:

\begin{itemize}
  \item USB, que en la placa da acceso a JTAG, para programar la FPGA y leer
    registros internos con SmartDebug, y a las consolas serie del MSS;
  \item Ethernet, en la misma red local, para operar la placa por SSH y para
    el protocolo de red.
\end{itemize}

La placa arranca Linux desde una tarjeta microSD. La computadora ejecuta
Libero SoC, las herramientas de verificación y el proceso de liberación.

\paragraph{Flujos de prueba}
Se usaron 39 flujos de conformidad MPEG-2 de distintos fabricantes y
laboratorios (CCETT, General Instrument, HHI, IBM, Nokia, Sony,
Tektronix, Toshiba, entre otros), con resoluciones de $32 \times 32$ a
$720 \times 576$. Para la demostración se usó un fragmento de \emph{Tears of
Steel} \citep{Blender_2012_tears}, codificado en MPEG-2 a $352 \times 288$ y
24 imágenes por segundo.

\paragraph{Referencia}
La referencia de corrección es \cmd{mpeg2decode}, el decodificador de
software del MPEG Software Simulation Group, incluido en el repositorio.

\section{Comparación de imágenes}
\label{sec:framecmp}

Saber si una imagen decodificada es correcta resultó más difícil de lo
esperado. La primera métrica daba un error absoluto medio cercano a 121 sobre
255 en imágenes que, vistas, eran correctas. La herramienta \cmd{framecmp}
resolvió las trampas que lo explicaban:

\begin{itemize}
  \item el núcleo guarda los píxeles como enteros con signo desplazados en
    $-128$, y se los estaba leyendo sin signo;
  \item en cada palabra de 64 bits, el píxel de la izquierda ocupa los bits
    más significativos, al revés que el orden de los bytes en memoria;
  \item la memoria de cuadros tiene cuatro buffers que rotan, y el orden de
    decodificación no es el de presentación: no hay una correspondencia fija
    entre buffer e imagen, así que cada buffer se compara contra todas las
    imágenes de referencia y se toma la mejor;
  \item los nombres internos de los planos de crominancia están cruzados
    respecto del estándar;
  \item un flujo cortado en cualquier byte deja una imagen parcial, y un
    volcado tomado durante la reconstrucción muestra una imagen a medio
    escribir: los flujos se cortan en frontera de imagen y el volcado espera
    a que la memoria deje de cambiar.
\end{itemize}

Para cada plano (Y, Cb, Cr) se reportan el error absoluto medio (MAD) y el
error pico, en niveles de 0 a 255. El método se validó primero en
simulación, hasta obtener cero donde debía haber cero, y recién después se
usó con el hardware.

\section{Regresión}
\label{sec:regresion}

El \emph{harness} \cmd{regress} ejecuta, desde la computadora y en unos 16
minutos, cuatro pruebas contra la placa:

\begin{description}[style=nextline, leftmargin=1.5em]
  \item[Conformidad] Cada flujo, cortado a 8 imágenes, se decodifica en la
    placa y se compara contra la referencia. Cada flujo recibe un veredicto
    según su peor buffer: \emph{bit-exacto} (MAD 0), \emph{cercano}
    (MAD $\leq 1$) o \emph{deriva} (MAD $> 1$), o \emph{sin escritura} si
    el núcleo no escribió ninguna imagen.
  \item[Determinismo] Un mismo flujo, cinco veces con \emph{reset} del
    núcleo y cinco cambiando de flujo sin \emph{reset}: todas las
    decodificaciones deben ser idénticas.
  \item[Robustez] Tres entradas dañadas (un flujo truncado, 200 bits
    invertidos en los datos y 256~KiB de bytes aleatorios), cada una seguida
    de un flujo correcto sin \emph{reset}, que debe decodificarse bit-exacto.
  \item[Velocidad] La memoria de cuadros se llena con un patrón conocido, se
    inicia el DMA y se mide el tiempo hasta la última escritura en la
    memoria de cuadros. Cinco corridas por flujo, después de una de
    calentamiento.
\end{description}

Los resultados se comparan contra una línea de base versionada, y la
comparación es insensible a la rotación de los buffers, que varía entre
corridas sin que cambie nada.

\section{Estimación de potencia}
\label{sec:modelo-potencia}

La Discovery Kit no permite medir el consumo por software: no tiene monitores
de corriente ni reguladores con interfaz digital, y sus puntos de prueba solo
sirven para medir tensión. Por eso la potencia se estimó con SmartPower, la
herramienta de Microchip, sobre la netlist ya ubicada y ruteada, en modo sin
vectores (con tasas de conmutación estimadas) y con las frecuencias reales
de cada reloj. Un script propio suma el resultado por instancia para obtener
la potencia de cada módulo del diseño.
